% ============================================================================
% LANGENTWURF LE-005: Arbeitsumgebung iPad
% Profil: S (Senioren 65+) | Tier: 2 (Standard)
% ============================================================================

\documentclass[11pt,a4paper]{article}
\usepackage[utf8]{inputenc}
\usepackage[T1]{fontenc}
\usepackage[ngerman]{babel}
\usepackage{geometry}
\usepackage{fancyhdr}
\usepackage{lastpage}
\usepackage{xcolor}
\usepackage{tcolorbox}
\usepackage{enumitem}
\usepackage{longtable}
\usepackage{hyperref}
\usepackage{amssymb}

\geometry{left=2.5cm, right=2.5cm, top=2.5cm, bottom=2.5cm}
\definecolor{fach}{HTML}{b35c00}
\definecolor{gray}{HTML}{666666}
\definecolor{successgreen}{HTML}{2a7a4b}

\tcbuselibrary{skins,breakable}
\newtcolorbox{scorebox}{ colback=white, colframe=fach, fonttitle=\bfseries, title=Didaktik-Score, breakable }

\pagestyle{fancy}
\fancyhf{}
\fancyhead[L]{\small\textcolor{gray}{Digitale Grundlagen -- 005 Arbeitsumgebung iPad}}
\fancyhead[R]{\small\textcolor{gray}{LE Tier 2 / Profil S}}
\fancyfoot[C]{\small\thepage{} / \pageref{LastPage}}
\renewcommand{\headrulewidth}{0.4pt}

\begin{document}

\begin{titlepage}
    \centering
    \vspace*{2cm}
    {\Large\textcolor{gray}{Digitale Grundlagen -- Senioren 65+}}\\[2cm]
    {\Huge\textcolor{fach}{\textbf{Langentwurf}}}\\[0.5cm]
    {\large Tier 2 | Profil S}\\[2cm]
    {\LARGE\textbf{005 -- Arbeitsumgebung iPad}}\\[0.5cm]
    {\large Modul A: Geräte \& Zugänge}\\[2cm]
    \begin{tabular}{rl}
        \textbf{Zeitrahmen:} & 60--75 Minuten \\
        \textbf{Vorkenntnisse:} & Einheit 004 (iPhone-Bedienung) \\
    \end{tabular}
    \vfill
    {\normalsize Stand: \today}
\end{titlepage}

\tableofcontents
\newpage

\section{Bedingungsanalyse}

\subsection{Ausgangssituation}
\begin{itemize}
    \item Nicht alle Teilnehmer haben ein iPad -- Einheit ist optional
    \item Wer iPhone bedienen kann, kann 80\% auch am iPad
    \item iPad-spezifische Features (Split View, Pencil) sind unbekannt
    \item Größerer Bildschirm bietet Vorteile für Senioren
\end{itemize}

\subsection{Typische Fragen}
\begin{itemize}
    \item ``Ist das iPad wie ein großes iPhone?''
    \item ``Kann ich am iPad telefonieren?''
    \item ``Brauche ich einen Stift dafür?''
\end{itemize}

\section{Sachanalyse}

\subsection{Kernkonzepte}
\begin{enumerate}
    \item \textbf{iPad = großes iPhone:} Gleiche Apps, gleiche Gesten, mehr Platz
    \item \textbf{Erweitertes Dock:} Mehr Apps, kürzlich verwendete rechts
    \item \textbf{Split View:} Zwei Apps nebeneinander (optional, Vertiefung)
    \item \textbf{Apple Pencil:} Stift zum Schreiben/Zeichnen (optional)
    \item \textbf{Bildschirmtastatur:} Größer, geteilte Tastatur möglich
\end{enumerate}

\subsection{Alltagsanalogien}
\begin{tabular}{|l|l|}
\hline
iPhone & Notizblock (kompakt, mobil) \\
iPad & Schreibtisch (mehr Platz, bequemer) \\
Split View & Zwei Dokumente nebeneinander auf dem Schreibtisch \\
Apple Pencil & Echter Stift für Notizen \\
\hline
\end{tabular}

\section{Didaktische Analyse}

\subsection{Stellung in der Reihe}
Optionale Vertiefung nach 004. Baut auf iPhone-Kenntnissen auf. Fokus auf \textbf{Unterschiede}, nicht auf Wiederholung.

\subsection{Didaktische Reduktion}
\textbf{Ausgeklammert:} Stage Manager, Slide Over, Tastatur-Shortcuts, iPad-spezifische Gesten mit 4 Fingern.

\textbf{Fokus:} Das iPad als ``komfortableres iPhone'' verstehen.

\section{Lernziele}

\begin{tabular}{|c|p{7cm}|c|c|}
\hline
\textbf{Nr.} & \textbf{Die Teilnehmer können...} & \textbf{Stufe} & \textbf{Niveau} \\
\hline
FZ1 & erklären, dass iPad und iPhone ähnlich funktionieren & Verstehen & Basis \\
FZ2 & das erweiterte Dock am iPad nutzen & Anwenden & Basis \\
FZ3 & Apps im Vollbild und in Split View öffnen & Anwenden & Vertiefung \\
FZ4 & den Apple Pencil für Notizen nutzen (falls vorhanden) & Anwenden & Vertiefung \\
\hline
\end{tabular}

\section{Verlaufsplanung}

\begin{longtable}{|p{1cm}|p{2cm}|p{5cm}|p{3cm}|p{2cm}|}
\hline
\textbf{Zeit} & \textbf{Phase} & \textbf{Inhalt} & \textbf{TN-Aktivität} & \textbf{Material} \\
\hline
0--10 & Einstieg & Vergleich: iPhone vs. iPad nebeneinander, Unterschiede benennen & Beobachten, Benennen & -- \\
\hline
10--25 & Gemeinsamkeiten & Gesten wiederholen, Apps öffnen, Kontrollzentrum & Mitmachen & S.1 \\
\hline
25--40 & Unterschiede & Dock (mehr Apps), größere Tastatur, Bildschirmgröße & Ausprobieren & S.1 \\
\hline
40--55 & Vertiefung & Split View zeigen (zwei Apps), Apple Pencil (falls vorhanden) & Optional mitmachen & S.2 \\
\hline
55--65 & Übungen & Checkliste, Vergleichstabelle ausfüllen & Bearbeiten & S.3 \\
\hline
\end{longtable}

\section{Lösungen}

\textbf{Was ist gleich?} Apps, Apple-ID, iCloud, Gesten, Einstellungen.

\textbf{Was ist anders?} Größerer Bildschirm, erweitertes Dock, Split View möglich, Pencil-Unterstützung, keine Telefonfunktion (außer mit WLAN-Anrufe).

\section{Didaktik-Score}

\begin{scorebox}
\begin{tabular}{|c|l|c|c|p{3.5cm}|}
\hline
\# & Dimension & Gew. & Score & Notiz \\
\hline
1 & Differenzierung & 15\% & 9/10 & Klar Basis/Vertiefung getrennt \\
2 & Taxonomie & 10\% & 9/10 & Transfer von iPhone-Wissen \\
3 & Alltagsrelevanz & 10\% & 9/10 & Großer Bildschirm = echte Erleichterung \\
4 & Aktivierung & 10\% & 9/10 & Viel Ausprobieren \\
5 & Scaffolding & 10\% & 9/10 & Baut auf 004 auf \\
6 & Feedback & 10\% & 9/10 & Vergleichstabelle als Selbstkontrolle \\
7 & Sprachliche Klarheit & 10\% & 9/10 & Vergleiche sind klar \\
8 & Motivation & 10\% & 9/10 & ``Wie iPhone, nur besser'' motiviert \\
9 & Barrierefreiheit & 10\% & 10/10 & iPad ist barrierefreundlicher \\
10 & Praktikabilität & 5\% & 9/10 & 65 Min angemessen \\
\hline
\multicolumn{3}{|r|}{\textbf{GESAMT:}} & \textbf{9.05/10} & \textcolor{successgreen}{\textbf{FREIGABE}} \\
\hline
\end{tabular}
\end{scorebox}

\end{document}
